\chapter*{Introduction générale}
\addcontentsline{toc}{chapter}{Introduction générale} % to include the introduction to the table of content
\markboth{Introduction générale}{} %To redefine the section page head

Le changement climatique est aujourd'hui l'un des défis majeurs auxquels fait face notre société. Une part importante des émissions de gaz à effet de serre provient directement des choix quotidiens des individus~: mode de transport, consommation d'énergie et d'eau, habitudes alimentaires, gestion des déchets ou encore usage des appareils électroniques. Or, la majorité des personnes n'ont aucun moyen simple, fiable et personnalisé d'estimer l'impact réel de leurs habitudes sur l'environnement. Les calculateurs d'empreinte carbone existants se limitent, dans leur grande majorité, à l'application de formules fixes et génériques, incapables de s'adapter à la diversité des profils d'usage et de fournir des recommandations réellement personnalisées.

C'est dans ce contexte qu'a été conçu et développé \textbf{GREENLY}, une plateforme web intelligente de calcul et de suivi de l'empreinte environnementale individuelle. Ce projet propose une approche hybride originale~: au lieu de reposer uniquement sur des formules statiques, GREENLY entraîne de véritables modèles de \textbf{Machine Learning} sur des données publiques réelles (agence américaine de protection de l'environnement EPA, jeux de données de l'université UCI, base \textit{Our World in Data}) afin de prédire, pour six domaines d'usage (transport, énergie, électronique, eau, alimentation, déchets), l'empreinte carbone ou hydrique associée. Au-delà de la simple prédiction, la plateforme intègre également un volet \textbf{Business Intelligence} complet~: tableaux de bord interactifs, indicateurs clés de performance, visualisations de tendances et module de prévision statistique, permettant à chaque utilisateur de suivre l'évolution de son impact environnemental dans le temps et d'identifier les leviers d'amélioration les plus efficaces.

L'objectif de ce projet est double. D'une part, il s'agit de démontrer la faisabilité et la pertinence d'une approche par apprentissage automatique pour l'estimation d'empreinte carbone, là où les outils existants se contentent de formules déterministes. D'autre part, il s'agit de concevoir une architecture logicielle complète, moderne et déployable~: une interface utilisateur réactive développée avec React et TypeScript, un micro-service de Machine Learning exposé via une API FastAPI, et une base de données Supabase (PostgreSQL) sécurisée, le tout orchestré via des conteneurs Docker.

Pour mener à bien ce travail, une démarche méthodologique inspirée du cadre Agile Scrum a été adoptée, permettant un développement itératif et incrémental, avec une identification progressive des besoins fonctionnels et non fonctionnels du système.

Ce rapport est structuré en quatre chapitres. Le premier chapitre présente le contexte général du projet, la problématique posée, les objectifs visés ainsi qu'une étude comparative de l'état de l'art des solutions existantes en matière de calcul d'empreinte carbone. Le deuxième chapitre est consacré à la spécification des besoins fonctionnels et non fonctionnels du système, à la modélisation des cas d'utilisation, ainsi qu'à la planification du projet selon la méthodologie Scrum (\textit{Product Backlog}, sprints, planning). Le troisième chapitre détaille la conception et la réalisation du volet Intelligence Artificielle~: architecture globale, diagrammes UML, modélisation de la base de données, préparation des jeux de données, entraînement et évaluation des modèles de Machine Learning, et intégration de ces modèles dans l'application. Le quatrième et dernier chapitre présente le volet Business Intelligence et prévision~: architecture décisionnelle, indicateurs clés de performance, tableaux de bord, visualisations, prévisions statistiques et recommandations générées automatiquement. Le rapport se termine par une conclusion générale qui dresse le bilan du travail réalisé et propose des perspectives d'évolution.
